\section{Private payment with public obligation}
\label{sec:private-payment}\label{sec:private-renewal}

Every node on the network must be able to decide, on its own and without asking anyone, whether the
workload it is about to run has been paid for. That decision has to be reproducible by thousands of
mutually distrusting machines from public data alone, which is precisely why the payment that answers
it is, today, a transparent transaction whose inputs name its sender. The tension this section
addresses is one sentence long: \emph{the network must be convinced that a deployment is funded while
learning neither who paid nor which payer belongs to which workload.} The workload itself is public
and must stay public --- a node cannot run a container whose image it is not told --- so the object to
hide is not the obligation but the payer, and the edge between the payer and the obligation. What
follows states that requirement formally, shows which parts of it the deployed Sapling toolkit already
satisfies, shows which part is unsatisfiable by construction and says so rather than omitting it,
measures the anonymity set that actually exists rather than assuming one, and names the single
consensus rule that gates the whole thing. The most important result in this section is a
\emph{negative scope} result: the general problem is smaller than it looks, and the part that remains
genuinely open is confined to one optional path.

\begin{remark}[Status of this section after the retirement decision]
\label{rem:pp-status}
\textbf{Flux is retiring both shielded pools} (\Cref{sec:shielded-retirement},
\srcintent{08-answers-round-2.md, round 10}). Every \emph{construction} in this section is therefore
inapplicable to the target architecture: there will be no pool to pay into, no Sapling output to
build a payment from, no anonymity set to grow, and no proving system to select. The section is
retained, unmodified in its technical content, for three reasons, and a reader should hold all three
in mind while reading it.

\begin{enumerate}
  \item \textbf{The requirement survives the decision.} \Cref{def:pp-pi} and \Cref{def:pp-ru} state what private
  payment for a public obligation must satisfy. That statement is independent of any construction and
  independent of whether Flux carries a shielded pool. If Flux ever revisits privacy --- via
  Orchard/Halo2 or otherwise --- this is the specification the work must meet, and it is written
  down here rather than needing to be rediscovered.
  \item \textbf{Two of its results are why retirement is defensible.}
  \S\ref{sec:pp-amount} proves amount hiding is unsatisfiable under \PNR{}, and
  \S\ref{sec:pp-scope} shows the residual problem is confined to one optional path. Together they
  establish that a shielded pool was buying considerably less than it appears to --- which is a load-
  bearing input to the decision, not a consolation after it.
  \item \textbf{The measurement is a permanent record.} \S\ref{sec:pp-anonymity} measures the
  anonymity set that actually existed on Flux. After retirement that measurement cannot be repeated,
  and it is the honest empirical answer to how much privacy the deployed pool ever provided.
\end{enumerate}

Read the constructions below as \emph{what would have been required}, not as what is planned. Where
the text says \planned{} of a shielded construction, that tag is now superseded by the retirement
decision and should be read as conditional on a pool existing.
\end{remark}

\begin{remark}[Symbols introduced here]
This section uses, in addition to \texttt{notation.tex}: $m_\app$ for the registration or update
message of application $\app$; $H_\app \defeq \mathsf{SHA256}(m_\app)$ for its hash, the value FluxOS
places in \texttt{OP\_RETURN}; $e_\app$ for the paid-through height; $A_{\mathrm{sink}}$ for the
payment sink, \S7's $\mathsf{sink}$; $v$ for a paid amount in \sat; $v_{\mathrm{bal}}$ for a Sapling \texttt{valueBalance};
$u$ for an individual FluxOS node; $\mathsf{AS}$ for an anonymity set; $\epsilon_{\mathrm{sc}}$ for a
side-channel advantage; $\tree_{\mathrm{ent}}$ and $\mathsf{rt}$ for an entitlement tree and its root;
$s$ for a credential secret; $j$ for an epoch index and $T_{\mathrm{ep}}$ for an epoch length; and
$\mathsf{xsk}$ for a ZIP-32 child spending key. These are proposed additions to the notation index
(Appendix~D) and are used nowhere else in the paper.
\end{remark}

\subsection{The verifier, and the requirement pair}
\label{sec:pp-requirements}

\subsubsection{What is public by construction}

Three things are public because the system cannot function otherwise, and it is worth separating them
before any privacy claim is made.

\begin{enumerate}
  \item \textbf{The workload.} An application's public part
    $\mathsf{spec}_{\mathrm{pub}}(\app) = (\mathsf{name}, H_\app, \mathsf{owner}, \ninst, \dur,
    \hgt_0, \mathsf{enterprise?})$ is stored by every node, because placement is opportunistic
    self-selection and any node may be the one that runs $\app$ (\shipped, \srccode{flux/\allowbreak{}ZelBack/\allowbreak{}src/\allowbreak{}services/\allowbreak{}appDatabase/\allowbreak{}registryManager.js}{1714}).
    Content is a separate matter: of the \num{1195} deployed applications, \num{972} publish their
    \texttt{compose} block in the clear and \num{223} (\qty{18.7}{\percent}) are enterprise records
    whose \texttt{compose} and \texttt{contacts} are stored zeroed
    \srcmeasured{api.runonflux.io/\allowbreak{}apps/\allowbreak{}globalappsspecifications}{2026-09-03}.
  \item \textbf{The paid-through height.} $e_\app$ must be a public function of the chain, or nodes
    disagree about whether to keep running a container.
  \item \textbf{The amount, under \PNR.} This is derived rather than chosen, and
    \S\ref{sec:pp-amount} proves it.
\end{enumerate}

What is \emph{not} required to be public is the identity of the party that moved the value, and the
association between that party and the application. That is the whole of the privacy goal here, and
stating it this narrowly is what makes the rest of the section tractable.

\subsubsection{The verifier}

\begin{definition}[Funding scheme and verifier]
\label{def:pp-scheme}
Let $\mathsf{Apps}$ be the set of application identifiers, $\mathcal{C}$ the set of chain prefixes,
$\mathsf{Tx}$ the set of transactions and $\mathsf{Msg}$ the set of FluxOS messages. A \emph{funding
scheme} is a pair $(\mathsf{Fund}, \mathsf{Verify})$ with
\[
  \mathsf{Fund} : \principal \times \mathsf{Apps} \longrightarrow \mathsf{Tx} \times (\mathsf{Msg} \cup \{\varepsilon\}),
  \qquad
  \mathsf{Verify}_u : \mathsf{Apps} \times \mathbb{N} \times \mathcal{C} \longrightarrow \mathbb{N} \cup \{\bot\},
\]
where $\mathsf{Fund}$ is run by a principal and $\mathsf{Verify}_u$ is run by FluxOS node $u$ on its
local chain prefix $C_{\le\hgt}$ and its local message store, returning a paid-through height $e_\app$
or $\bot$. Funded status is $\mathsf{Funded}(\app,\hgt) \defeq \indic{\hgt \le e_\app}$.
\end{definition}

\begin{assumption}[Verifier model, from \S2]
\label{ass:pp-verifier}
Node $u$ holds $C_{\le\hgt}$, the FluxOS temporary and permanent message sets, and \emph{no secret}:
no shared key, no issuer key, and no viewing key it did not derive from public data. Node $u$ needs no
message from any other verifier to evaluate $\mathsf{Verify}_u$. This is a description of the deployed
system, not a design preference: \texttt{explorerService.\allowbreak{}processInsight} recognises a payment by output
address and \texttt{OP\_RETURN} and \texttt{messageVerifier.\allowbreak{}checkAndRequestApp} promotes the
registration message, on every node, from chain data alone (\shipped,
\srccode{flux/\allowbreak{}ZelBack/\allowbreak{}src/\allowbreak{}services/\allowbreak{}explorerService.js}{309--345},
\srccode{flux/\allowbreak{}ZelBack/\allowbreak{}src/\allowbreak{}services/\allowbreak{}appMessaging/\allowbreak{}messageVerifier.js}{619--897}); CumulusVPN states the same property
of its gateways --- ``every gateway runs the same deterministic scan and therefore reaches the same
\texttt{paid\_until} map [\ldots] with no server-to-server coordination'' (\shipped, \srccode{cumulusvpn/\allowbreak{}gateway/\allowbreak{}internal/\allowbreak{}entitle/\allowbreak{}entitle.go}{1--5}).
\end{assumption}

\begin{assumption}[Adversary, from \S2]
\label{ass:pp-adversary}
The adversary $\adv$ is a node operator. It sees the whole chain, the public specification of every
application, all FluxOS gossip, and the client traffic arriving at the nodes it runs; it may run any
number of nodes (the measured maximum under one signing key is \num{237}, \qty{3.7}{\percent} of the
fleet \srcmeasured{api.runonflux.io/\allowbreak{}daemon/\allowbreak{}listzelnodes}{2026-09-03}); it may be the \emph{owner} of
the application whose payer it wants to identify; and it may hold shielded funds. It is bounded by
$\advbudget$ for anything requiring capital and is otherwise probabilistic polynomial time.
\end{assumption}

\subsubsection{The hard requirements}

\begin{definition}[Public obligation]
\label{def:pp-obligation}
A funding scheme satisfies \emph{public obligation} if:
\begin{description}
  \item[V0 (no secret).] $\mathsf{Verify}_u$ is a deterministic function of $C_{\le\hgt}$ and public
    messages, for every $u$.
  \item[S (soundness).] For every PPT payer $\adv$,
    \[
      \Pr\!\left[\mathsf{Verify}_u(\app,\hgt,C_{\le\hgt}) = e > \hgt
      \ \wedge\ \textstyle\sum \text{value confirmed to } A_{\mathrm{sink}} \text{ bound to } H_\app \text{ by } \hgt < \price(\app)\right]
      \le \mathrm{negl}.
    \]
  \item[C (consistency).] Honest $u, u'$ with the same chain prefix return the same $e_\app$.
  \item[L (liveness under partition).] $\mathsf{Verify}_u$ requires no message from any other party; a
    partitioned node evaluates on its local prefix and converges when the prefix does.
  \item[D (no double-funding).] One confirmed transfer contributes to at most one pair
    $(\app,\text{period})$.
\end{description}
\end{definition}

\begin{definition}[Payer indistinguishability, PI]
\label{def:pp-pi}
$\adv$ names two principals $\principal_0,\principal_1$, each able to fund, and an application $\app$;
$\adv$ may own $\app$ and any set of nodes. The challenger draws $b \in \{0,1\}$ uniformly and runs
$\mathsf{Fund}(\principal_b,\app)$; $\adv$ observes the resulting chain and network state and outputs
$b'$. Define
\[
  \mathrm{Adv}^{\mathrm{PI}}(\adv) \defeq \left|\Pr[b'=b] - \tfrac12\right| \in [0,\tfrac12].
\]
The scheme has \emph{cryptographic} PI if $\mathrm{Adv}^{\mathrm{PI}}(\adv) \le \mathrm{negl}$ whenever
$\principal_0$ and $\principal_1$ fund from the same pool, at the same time, with the same amount. Its
\emph{effective} advantage is $\mathrm{negl} + \epsilon_{\mathrm{sc}}$, where $\epsilon_{\mathrm{sc}}$
collects timing, amount and network-metadata leakage. \S\ref{sec:pp-anonymity} measures
$\epsilon_{\mathrm{sc}}$'s inputs rather than assuming them away.
\end{definition}

\begin{definition}[Renewal unlinkability, RU]
\label{def:pp-ru}
\emph{Payer-side:} for two funding events of the same $\app$ at periods $i \ne j$, $\adv$ cannot
distinguish whether one principal produced both. \emph{Cross-application:} for $\app \ne \app'$,
$\adv$ cannot distinguish whether one principal funded both. Two renewals of the same $\app$ are
linked by $H_\app$ by construction --- that is the obligation being public --- so RU never denotes
hiding \emph{that} $\app$ was renewed.
\end{definition}

\begin{definition}[Amount hiding, AH]
\label{def:pp-ah}
$\adv$ learns nothing about $v$ beyond $v \ge \price(\app)$.
\end{definition}

\begin{definition}[Forward secrecy, FS]
\label{def:pp-fs}
Compromise at period $i$ of all key material used to produce funding events of periods $\ge i$ does
not increase $\mathrm{Adv}^{\mathrm{PI}}$ and does not break RU for periods $< i$.
\end{definition}

Definitions~\ref{def:pp-obligation} and~\ref{def:pp-pi} are the \emph{hard} requirements;
Definitions~\ref{def:pp-ru}--\ref{def:pp-fs} are desirable, and \S\ref{sec:pp-shapes} states which are
achievable in which shape. Content privacy --- the hiding of $\mathsf{spec}_{\mathrm{priv}}(\app)$
from non-selected nodes by enterprise encryption --- is owned by \S10 and constrained here only in
that a funding scheme must not reduce it beyond what $\price(\app)$ already reveals. AH is the sole
point at which the two interact, which brings us to the first result.

\subsection{Amount hiding is unsatisfiable under PNR}\label{sec:pp-amounthiding}
\label{sec:pp-amount}

This is stated early because it is a genuine conflict between two of this paper's own mechanisms, and
because a reader who expects Zcash-style amount privacy on hosting payments should be disabused before
reading a construction that does not deliver it.

\begin{property}[Amount recovery from the fee pool]
\label{prop:pp-amount}
Let $\nodeshare(\hgt) \in (0,1]$ be the operator share at height $\hgt$ and write
$r(\hgt) = 100\,\nodeshare(\hgt) \in \{1,\dots,100\}$. Under \PNR{} (\S7, PT1--PT5), a hosting payment
of price $\price \in \mathbb{N}$ \sat{} credits
$A_\Phi = \left\lfloor r(\hgt)\,\price / 100 \right\rfloor$ \sat{} to the consensus fee pool $\pool$,
and the per-block drip $D_\hgt = \floorfrac{\pool_{\hgt-1}}{\drip}$ appears in the coinbase. Then
$A_\Phi$ is a public function of the chain, and consequently
\[
  \price \in \left[\ \frac{100\,A_\Phi}{r(\hgt)},\ \frac{100\,(A_\Phi+1)}{r(\hgt)} \right),
  \qquad \text{an interval of width } \frac{100}{r(\hgt)} \sat.
\]
At $\nodeshare = 0.20$ the price is recoverable to within \num{5}\sat; at the $\nodesharecap = 0.50$
ceiling, to within \num{2}\sat. Therefore \textbf{no funding scheme that settles through \PNR{}
satisfies AH} (Definition~\ref{def:pp-ah}).
\end{property}

\begin{proof}
Every node must compute $\pool_\hgt$ exactly in order to validate the coinbase, since $D_\hgt$ is a
consensus-checked function of $\pool_{\hgt-1}$ and a node that cannot compute the pool cannot decide
block validity. Hence $\pool_\hgt$ is a deterministic public function of $C_{\le\hgt}$, and so is each receipt's
$A_\Phi$: under PS4 of \S7 it is computed from the value of a transparent output paying
$A_{\mathrm{sink}}$, and the block's inflow $F_\hgt$ is the sum of these. Inverting the floor gives the stated
interval, whose width is $100/r$. Any scheme achieving AH would have to hide $A_\Phi$ from at least
one honest verifier, contradicting checkability of the coinbase; a hidden pool cannot drip a checkable
coinbase.
\end{proof}

\begin{remark}[\S7 and \S21 cannot both have what they want]
\label{rem:pp-conflict}
The resolution recorded in this paper is that amount privacy is \textbf{out of scope for
\PNR-typed payments, by construction rather than by omission}. The cost is asymmetric: for the
\num{972} public-content applications it is nil, because $\price$ is already a public function of a
public specification; for the \num{223} enterprise applications the price reveals
$\langle \pricevec, \resvec \rangle$, a coarse size fingerprint of encrypted content. The only escape
--- exempting a payment class from \PNR{} --- contradicts \S7's requirement that all hosting revenue
settle through the split, and is left as an open decision (Q5, Table~\ref{tab:pp-open}).
\end{remark}

\subsection{Three shapes, kept apart}
\label{sec:pp-shapes}

Most of the difficulty attributed to ``private payment on Flux'' comes from conflating three problems
whose hardness differs sharply. Table~\ref{tab:pp-shapes} separates them, and every later claim in
this section is scoped to one of the three.

\begin{table}[H]
\centering
\small
\caption{The three shapes of the private-payment problem. Only shapes B and C carry a standing object
that must be re-proved; shape A does not, which is the content of
Proposition~\ref{prop:pp-scope}.}
\label{tab:pp-shapes}
\begin{tabular}{@{}p{0.06\linewidth}p{0.20\linewidth}p{0.24\linewidth}p{0.20\linewidth}p{0.18\linewidth}@{}}
\toprule
& \textbf{The obligation names} & \textbf{Example} & \textbf{What must be hidden} & \textbf{Maturity} \\
\midrule
\textbf{A} & a public workload $H_\app$, which every node must know in order to run it &
application registration and renewal; the direct path of \S16 &
the payer &
\planned; deployed circuits, gated on \S\ref{sec:pp-blocker} \\
\addlinespace
\textbf{B} & a private user: a pseudonym whose only public attribute is ``paid until $e$'' &
CumulusVPN, memo $=\mathsf{base58}(\mathsf{SHA256}(pk_t)[0{:}20])$ &
the link pseudonym\,$\leftrightarrow$\,payment, \emph{and} pseudonym\,$\leftrightarrow$\,payer &
\vision; needs a circuit Flux does not have \\
\addlinespace
\textbf{C} & a standing balance debited over time &
deposited credits with auto-renewal, the credit path of \S16 &
the payer, and the link between successive debits &
\planned; reduces to A when the renewer is the payer's own agent \\
\bottomrule
\end{tabular}
\end{table}

Table~\ref{tab:pp-props} then states, per shape, which requirement is hard, which is desirable, and
which cannot be had. It is the map for the rest of the section.

\begin{table}[H]
\centering
\small
\caption{Requirements per shape. ``hard'' = a mechanism is not acceptable without it; ``desirable'' =
wanted, and its absence is stated rather than hidden.}
\label{tab:pp-props}
\begin{tabular}{@{}lllll@{}}
\toprule
\textbf{Property} & \textbf{A, direct} & \textbf{C, own agent} & \textbf{C, third-party renewer} & \textbf{B, CumulusVPN-class} \\
\midrule
V0, S, C, L & hard & hard & hard & hard \\
D & hard & hard & hard & vacuous (time-based) \\
PI vs.\ the network & hard & hard & hard & hard, vs.\ the gateway \\
PI vs.\ the renewer & n/a & n/a & desirable; \textbf{open} (R2) & n/a \\
PI, transparent payer & impossible & --- & --- & desirable; needs A \\
RU payer-side & from PI per event & hard & hard & hard, across epochs \\
RU cross-application & desirable & desirable & desirable & --- \\
AH & \textbf{unsatisfiable} (Prop.~\ref{prop:pp-amount}) & unsatisfiable & unsatisfiable & n/a, fixed price \\
FS & desirable & desirable & desirable & desirable \\
\bottomrule
\end{tabular}
\end{table}

\subsection{The scope result}
\label{sec:pp-scope}

The proposition below is the most consequential finding in this section, because it decides how much
of the paper's agent-native argument depends on an unsolved problem. It says: not much.

\begin{proposition}[No new relation is required for shape A]
\label{prop:pp-scope}
Let $\mathsf{Verify}_u$ be the deployed FluxOS funding rule of \S\ref{sec:pp-deployed}. Then:
\begin{enumerate}
  \item[(i)] $\mathsf{Verify}_u(\app,\hgt,C_{\le\hgt})$ depends on $C_{\le\hgt}$ only through the set
    of confirmed \emph{funding events}, where a funding event is a confirmed transaction carrying a
    transparent output $(A_{\mathrm{sink}}, v)$ together with $\mathsf{OP\_RETURN} = H_\app$. Between
    funding events, no predicate over hidden state is evaluated.
  \item[(ii)] There exists a funding scheme satisfying V0, S, C, L, D and cryptographic PI
    (Definitions~\ref{def:pp-obligation}, \ref{def:pp-pi}) whose only zero-knowledge component is the
    Sapling spend statement and the Sapling output statement, both already verified by every
    \fluxd{} node at consensus (\shipped,
    \srccode{fluxd/\allowbreak{}src/\allowbreak{}main.cpp}{1447--1482}, \srccode{fluxd/\allowbreak{}src/\allowbreak{}primitives/\allowbreak{}transaction.h}{155--242}).
  \item[(iii)] Consequently no relation outside the deployed circuits is required for shape A, and in
    particular ``a statement proving that a note was spent to a known application address without
    revealing the sender'' is not a new circuit: it is the Sapling spend statement with a transparent
    recipient.
\end{enumerate}
\end{proposition}

\begin{proof}[Proof sketch]
\emph{(i)} By inspection of the deployed rule (\S\ref{sec:pp-deployed}): $e_\app$ is assigned once per
funding event from $(v, \hgt', \dur(m_\app))$ and is never re-derived; $\mathsf{Verify}_u$ thereafter
compares two integers. There is therefore no standing obligation and nothing to prove repeatedly.
\emph{(ii)} Take Construction~\ref{con:pp-a} below. Its verifier is unchanged from the deployed one
because the FluxOS scanner keys on outputs and \texttt{OP\_RETURN} and never on inputs; V0, S, C, L, D
hold by the same argument as for the transparent path, since the value transfer is still a confirmed
transparent output to $A_{\mathrm{sink}}$. Cryptographic PI reduces to Sapling spend
indistinguishability: two funding transactions produced by $\principal_0$ and $\principal_1$ spending
notes of sufficient value differ only in $\nf$, $\mathsf{cv}$, $\mathsf{rk}$, $\zkproof$ and the
change-note ciphertext, each of which is hiding under Sapling's stated assumptions
\citep{zcashprotocol}; the number of spend and output descriptions, the anchor and the fee are not
hidden, so the builder must hold them equal across payers or they fall under $\epsilon_{\mathrm{sc}}$.
\emph{(iii)} Immediate from (i) and (ii): the verifier needs $v \ge \price$,
the binding of the transfer to $H_\app$, and that $v$ actually moved. The first two are transparent
fields; the third is exactly what the binding signature and the spend statement establish. No
statement \emph{about the application} is proved in zero knowledge, because the application is public.
\end{proof}

\begin{corollary}[Agent-native direct payment]
\label{cor:pp-agent}
When the renewing party is the tenant's own agent, shape C reduces to shape A composed with ZIP-32
child-key derivation (\shipped, \srccode{fluxd/\allowbreak{}src/\allowbreak{}flux/\allowbreak{}zip32.h}{53--136}): the ``balance'' is the
value of the notes held under a derived child key, it exists nowhere as protocol state, and each
renewal is an independent shape-A funding event. \textbf{Agent-native direct payment therefore
collapses to a plain Sapling spend plus ZIP-32.}
\end{corollary}

\begin{corollary}[Blast radius]
\label{cor:pp-blast}
The hard version of the problem --- repeated unlinkable proofs against a hidden standing object ---
arises only when something hidden must be \emph{reused} by a party that must not learn the link. That
is (a) shape C with a \emph{third-party} renewer, and (b) shape B for payers who cannot or will not
use the shielded pool. The agent-native argument of \S18 survives with the general problem still open,
and the exposure is limited to the optional credit path of \S16.
\end{corollary}

\begin{remark}[Constructible is not private]
\label{rem:pp-constructible}
Proposition~\ref{prop:pp-scope} is a statement about the \emph{relation to be proved}, not about
deployability or about realised privacy. Construction~\ref{con:pp-a} is constructible over deployed
circuits; whether it is private in practice depends on the size of the payer's anonymity set
(\S\ref{sec:pp-anonymity}) and on a consensus rule that currently forbids entering the pool at all
(\S\ref{sec:pp-blocker}). The paper claims only the first until the measurement M1 is taken. The
limitation is a result, not a consolation: knowing that three-quarters of the problem needs no new
cryptography is worth more than a construction for all of it that nobody has.
\end{remark}

\subsection{The blocker: the shielded pool cannot be entered}
\label{sec:pp-blocker}

\begin{assumption}[The pool is closed --- \shipped, and it gates everything below]
\label{ass:pp-closed}
Since \texttt{UPGRADE\_FLUX} at mainnet height \num{835554},
\texttt{Contextual\allowbreak{}Check\allowbreak{}Transaction} rejects any transaction that carries transparent inputs
($\mathtt{tx.vin.size()} > 0$) together with any new shielded output --- any non-empty
\texttt{vJoinSplit} or \texttt{v\allowbreak{}Shielded\allowbreak{}Output} --- with reject reason
\texttt{bad-\allowbreak{}txns-\allowbreak{}fluxrebrand-\allowbreak{}vprivacy-\allowbreak{}not-\allowbreak{}empty}
(\shipped, \srccode{fluxd/\allowbreak{}src/\allowbreak{}main.cpp}{1398--1408}). The in-source comment states the intent: ``we no
longer want to allow coins to go into the privacy pool''. The rule does not inspect
\texttt{tx.\allowbreak{}v\allowbreak{}Shielded\allowbreak{}Spend}, so $\zaddr \to \zaddr$ and $\zaddr \to \taddr$ transactions remain valid.
\textbf{Shielded value can move and can leave; it cannot enter.} This is a Flux-original consensus
policy, not an inherited Zcash ZIP --- a distinction \S20 must preserve.
\end{assumption}

\begin{corollary}[Who can use construction A today]
\label{cor:pp-who}
Construction~\ref{con:pp-a} would be available only to holders of value shielded before
Assumption~\ref{ass:pp-closed} took effect in April 2021: \num{77188.75921324}\flux{} in the Sapling
pool and \num{19291.17735041}\flux{} in the Sprout pool at chain tip \num{2917115}, totalling
\num{96479.93656365}\flux{}
\srcmeasured{api.runonflux.io/\allowbreak{}daemon/\allowbreak{}getblockchaininfo (valuePools)}{2026-09-03}. Against issued
supply of \num{429466823.97}\flux{} (\S5) that is \qty{0.0225}{\percent}. The pool has been closed for
\num{2081561} blocks, approximately \num{5.4} years of wall clock, and its value is
monotonically non-increasing: every $\zaddr \to \taddr$ hosting payment removes a participant's value
from it and can never be replaced.
\end{corollary}

Two consequences deserve to be stated plainly rather than buried. First, the roadmap's privacy items
--- ``we are bringing private transactions back'', ``shielded by default in the wallets we make'' ---
were \textbf{never wallet-integration tasks}: re-opening $\taddr \to \zaddr$ is a consensus change.
That distinction outlives the decision that followed it, because it is what made the roadmap's
sequencing wrong rather than merely optimistic. \textbf{Reopening is now cancelled}: both pools are
to be retired instead \srcintent{08-answers-round-2.md, round 10} (\planned,
\S\ref{sec:shielded-retirement}). Second, capacity was never the constraint and membership was: at
the list-price revenue of the entire measured population --- \num{3817.62}\flux{} per
\num{30.56}-day period under \Cref{def:price} as corrected, of which \num{1784.97}\flux{}
(\SI{46.8}{\percent}, over \num{247} applications) is deployed by the Foundation's own identities and
\num{2032.65}\flux{} (\SI{53.2}{\percent}, over \num{948}) is third-party
\srcdoc{scripts/21-private-renewal.py} --- the Sapling pool alone would have funded the whole
network's hosting for \num{1.69} years, or \num{3.18} years of third-party hosting alone.

The split is taken on the specification's \texttt{owner} field --- the ZelID --- and not on the
address that paid, and the distinction is load-bearing rather than pedantic. Under the custodial path
of \S\ref{sec:paying-for-capacity} the service holds both the user's ZelID key and their payment key
and signs on their behalf (C42), so a card-paid customer deployment settles \emph{from a
Foundation-held address while belonging to the customer's ZelID}. Attributing by payer would count
every such deployment as the Foundation's own and understate third-party demand by however much of it
arrives that way. \settlednote{attribution is by \texttt{owner} ZelID. Of the four
Foundation-configured identities, three own nothing or nearly nothing --- \texttt{usersToExtend}
(current) and \texttt{fluxTeamFluxID} own no applications and \texttt{fluxSupportTeamFluxID} owns one
--- which is what the configuration implies, since \texttt{usersToExtend} confers permission to
extend an application \emph{on behalf of its owner} rather than ownership
\srccode{flux/\allowbreak{}ZelBack/\allowbreak{}config/\allowbreak{}default.js}{229}. The entire
Foundation share therefore rests on one address, \texttt{196GJWyLxzAw\ldots}, holding 246 of the
\num{1195} applications. Its composition supports the attribution --- \texttt{FluxWhitepaper},
\texttt{FluxWordPress} and \texttt{FluxRosettaServer} alongside \texttt{BitcoinWhitepaper},
\texttt{Aave}, \texttt{Electrum}, \texttt{EthereumNodeLight}, \texttt{Firod} and the 26 folding
deployments --- and the author confirms it as the team ZelID
\srcintent{08-answers-round-2.md, round 21}. The configuration itself names that address only as a
former extension-permission holder, so the identification rests on that confirmation rather than on
any declaration in code.} A pool
used for hosting by a handful of pre-2021 holders identifies them as the only possible payers, which
is a smaller anonymity set than the transparent path offers.

That second point is what survives retirement, and it deserves emphasis because it bears directly on
how much was given up. \Cref{cor:pp-who} shows the anonymity set of the deployed pool was bounded by
who held shielded value before 2021 --- a set that could not grow while $\taddr \to \zaddr$ stayed
closed, and that reopening would have grown only slowly and only for new entrants. The practical
privacy on offer at the moment the decision was taken was therefore thin: not zero, but a set small
enough that using it to pay for hosting would in many cases have been self-identifying. A reader
weighing this decision should weigh it against that measured set, not against the privacy a mature
shielded pool would provide.

\begin{remark}[Q1 and Q2 --- both resolved by retirement, in opposite directions]
\label{rem:q1q2-resolved}
Q1 asked \emph{when} $\taddr \to \zaddr$ would reopen, having been narrowed from \emph{whether}. Q2
asked whether ZIP-209 turnstile enforcement --- compiled in but dormant against unmodified
upstream-Zcash checkpoint constants --- should be activated, and in what order relative to Q1.

\textbf{Q1 is withdrawn.} There is no activation height to choose, because there is no reopening.

\textbf{Q2 is answered yes, and becomes more urgent rather than less.}
\S\ref{sec:shielded-retirement} recommends activating turnstile enforcement with Flux-correct
constants for the duration of the drain-only window. In that window the turnstile's check is exactly
the invariant that must hold --- the pool balance may only decrease --- so it is both cheap and
precisely correct, and it is the only period in which the pools shrink under a published deadline
with no offsetting inflow. Leaving enforcement dormant through the one window in which it is
trivially correct would be the worst of the available choices.
\end{remark}

\proposednote{Live for the drain-only window and moot after it (\Cref{rem:pp-status}, \Cref{rem:q1q2-resolved}). Q2. ZIP-209 turnstile enforcement, compiled in and dormant with unmodified upstream-Zcash
checkpoint constants (\srccode{fluxd/\allowbreak{}src/\allowbreak{}chainparams.cpp}{322--327}, \srccode{fluxd/\allowbreak{}src/\allowbreak{}chainparams.h}{200}).
Domain: \{on with Flux-computed checkpoints; off\}. Recommended: on, for the duration of the window;
round 11 confirmed the window's heights, not this switch \srcintent{08-answers-round-2.md, round 18}.
Trade-off --- on: a soundness failure in the
inherited 2018 parameters is bounded to the pool's declared value, and \S5's supply-verifiability
claim becomes true; off: a hosting payment could be printed from a forged proof undetectably. \S20
makes the same recommendation.}

\subsection{Why the naive construction fails: CumulusVPN, worked through}
\label{sec:pp-naive}

CumulusVPN is the shipped instance of the problem in this corpus, and it fails in the exact way the
abstract statement predicts. Its flow is (all \shipped):

\begin{enumerate}
  \item The client generates a WireGuard keypair $(sk_t, pk_t)$ on-device.
  \item \texttt{POST /v1/enroll \{pubkey, pow\_nonce\}} returns
    \texttt{\{payment\_address,\allowbreak{} payment\_memo,\allowbreak{} price\_flux,\allowbreak{} \dots\}}
    (\srccode{cumulusvpn/\allowbreak{}gateway/\allowbreak{}internal/\allowbreak{}api/\allowbreak{}api.go}{177,185--198}). \emph{The gateway now holds
    $pk_t$ and the client's IP address.}
  \item The client pays \texttt{price\_flux} --- \num{20}\flux{} per month as documented --- to
    \texttt{payment\_address} with \texttt{OP\_RETURN} memo \texttt{CVPN1:code}, where
    $\mathsf{code} = \mathsf{base58}(\mathsf{SHA256}(pk_t)[0{:}20])$, computed byte-identically in
    gateway and client (\srccode{cumulusvpn/\allowbreak{}gateway/\allowbreak{}internal/\allowbreak{}entitle/\allowbreak{}entitle.go}{99--108},
    \srccode{cumulusvpn/\allowbreak{}clients/\allowbreak{}core-ts/\allowbreak{}src/\allowbreak{}paymentCode.ts}{17--23}).
  \item Every gateway scans the chain independently; \texttt{ValidPayment} requires
    $30(v + 10^{-9}) \ge \mathsf{price}$ and exactly one \texttt{CVPN1:} memo, and grants
    $\lfloor 30 v / \mathsf{price}\rfloor$ days, capping \texttt{paid\_until[code]} at \num{720} days
    (\srccode{cumulusvpn/\allowbreak{}gateway/\allowbreak{}internal/\allowbreak{}entitle/\allowbreak{}entitle.go}{218--287}).
\end{enumerate}

As a public-obligation scheme this is correct and complete: V0 holds (chain only), S holds (value to
the address with the memo), C holds (identical derived maps), L holds (local scan), and D holds (the
memo resolves to one code; stacking is additive). It is the privacy half that fails, and it fails
totally.

\begin{property}[Enrolment collapses the payer anonymity set to one]
\label{prop:pp-cvpn}
Let $\adv$ be a gateway with which the client has enrolled, and let $\principal_0,\principal_1$ pay
from distinguishable transparent input sets. Then $\mathrm{Adv}^{\mathrm{PI}}(\adv) = \tfrac12$; that
is, $\adv$ wins the game of Definition~\ref{def:pp-pi} with probability $1$.
\end{property}

\begin{proof}
$\mathsf{code}$ is a deterministic function of $pk_t$ evaluated by the same code path in the client
and in the gateway's scanner. $\adv$ received $pk_t$ at step~2, so it computes $\mathsf{code}$; the
chain index is public and \texttt{ValidPayment} requires exactly one \texttt{CVPN1:} memo per
transaction, so the lookup returns a unique confirmed transaction; that transaction's transparent
inputs are public and identify the payer's wallet cluster under the ordinary common-input-ownership
and change-detection heuristics. $\adv$ therefore recovers $b$ with certainty.
\end{proof}

The one-way hash protects exactly the parties who never see its preimage. In a service whose verifier
must see the preimage in order to serve, it protects nobody who matters. What the gateway holds per
enrolled user is, concretely: $pk_t$; the client IP; last-handshake times, persisted to disk at
\texttt{/\allowbreak{}data/\allowbreak{}peers.\allowbreak{}cache} (\srccode{cumulusvpn/\allowbreak{}gateway/\allowbreak{}internal/\allowbreak{}wg/\allowbreak{}peercache.go}{47--51});
destination IPs and traffic timing; the code; and the paying address. Multi-hop does not help in v1,
because the same key is used on both hops, so the entry node computes the same code.

\begin{property}[Generalisation: it is the co-location, not the hash]
\label{prop:pp-general}
Every naive construction publishes one object with two sides: an \emph{identifying input side}
(transparent spends, which name a wallet cluster) and an \emph{identifying obligation side} (an
address or memo that names the workload or the user). The link between them is the transaction itself,
and it is public because the obligation must be checkable by parties holding no secret (V0). The leak
is therefore not the hash function, not the address format and not the memo encoding; it is the
co-location of the two sides in one public object. Exactly one of the two sides must be made
non-identifying, and the obligation side cannot be omitted.
\end{property}

Flux's own deployment path is the same object with different labels: transparent inputs, an output to
$A_{\mathrm{sink}}$ = \texttt{t3Nryf\allowbreak{}AQLGe\allowbreak{}Fs9j\allowbreak{}Eoeqsxm\allowbreak{}BN2QLRa\allowbreak{}RKFLUX}
\srcmeasured{api.runonflux.io/\allowbreak{}apps/\allowbreak{}deploymentinformation}{2026-09-03}, an \texttt{OP\_RETURN} carrying
$H_\app$, and a registration message signed by the owner key. Payer cluster, owner identity and
workload are all public and all linked.

\begin{property}[Payment-key indirection alone does not remove the leak]
\label{prop:pp-v15}
CumulusVPN's planned v1.5 separates a payment key $pk_p$ from the tunnel key $pk_t$
(\planned, no code found, \srcdoc{cumulusvpn/docs/04-payments.md:81--83}). This defeats an adversary
that sees only $pk_t$ --- a multi-hop entry node, a passive network observer --- because it cannot
compute $\mathsf{H}(pk_p)$. But the gateway must still decide whether $pk_t$ is entitled, and the only
way for the client to demonstrate that is to present $pk_p$ or a signature under it, which returns the
preimage to the gateway. Indirection therefore moves the linkage from enrolment to the entitlement
check; it does not remove it. It removes it only when the entitlement check is \emph{itself}
non-identifying, which requires either Construction~\ref{con:pp-b} or an issuer-signed voucher.
\end{property}

The paper states this so that ``v1.5 closes the gap'' is not printed anywhere. The two remedies
Property~\ref{prop:pp-general} permits are developed in \S\ref{sec:pp-consta} (blind the input side)
and \S\ref{sec:pp-constb} (blind the obligation side).

\subsection{The deployed rule, and construction A}
\label{sec:pp-deployed}

\subsubsection{What is deployed}

Per node, FluxOS holds the chain $C_{\le\hgt}$, a message store with temporary messages (1\,h TTL) and
permanent messages, and, per application name, a stored specification and $e_\app$. On observing a
funding event at height $\hgt'$, node $u$ runs \texttt{check\allowbreak{}And\allowbreak{}Request\allowbreak{}App}: it fetches $m_\app$ from
local temporary storage or from peers (up to three attempts), re-verifies the owner signature,
recomputes $\price(m_\app,\hgt')$ from the height-scheduled price table, and sets
\[
  e_\app \leftarrow
  \begin{cases}
    \hgt' + \dur(m_\app), & v \ge \price(m_\app,\hgt'), \\[2pt]
    \text{unchanged, warning logged}, & \text{otherwise,}
  \end{cases}
  \qquad
  \mathsf{Verify}_u(\app,\hgt,C_{\le\hgt}) = \indic{\hgt \le e_\app}.
\]
The second branch is the silent-underpayment defect of \S16: an underpaying transaction is consumed,
not refunded, and no signed rejection reaches the payer (\shipped,
\srccode{flux/\allowbreak{}ZelBack/\allowbreak{}src/\allowbreak{}services/\allowbreak{}appMessaging/\allowbreak{}messageVerifier.js}{619--897}). A shielded payer has no return
address, so the negative acknowledgement \S16 requires must be published on the FluxOS message layer
keyed by $H_\app$ for the payer to poll --- an interface requirement (I4 below), not a new mechanism.

From a transparent funding event a node learns everything: the inputs and hence the payer's address
cluster; $v$; $H_\app$; $m_\app$ in full; and, if it is the API node through which the registration
was submitted, the client's IP.

\subsubsection{Construction A}
\label{sec:pp-consta}

\begin{construction}[Shielded direct payment over the deployed toolkit --- \planned]
\label{con:pp-a}
A funding transaction for $\app$ would be a Sapling v4 transaction
(\texttt{SAPLING\_TX\_VERSION}~$=4$, \texttt{SAPLING\_VERSION\_GROUP\_ID}~$= \mathtt{0x892F2085}$)
with:
\begin{itemize}
  \item $\mathsf{vin} = \varnothing$;
  \item one or more \texttt{Spend\allowbreak{}Description}s spending the payer's notes;
  \item one \texttt{Output\allowbreak{}Description} returning change to a payer-controlled \zaddr;
  \item transparent outputs $(A_{\mathrm{sink}}, v)$ and $(\mathsf{OP\_RETURN}\ H_\app, 0)$;
  \item $v_{\mathrm{bal}} = v + \mathsf{fee}$, and a \texttt{bindingSig} over the value commitments.
\end{itemize}
Under \PNR{} ($\hgt \ge \hpa$) the same transaction would be unchanged in shape: its single sink
output $(A_{\mathrm{sink}}, v)$ is divided by consensus under PS4 of \S7, and $H_\app$ stays in the
\texttt{OP\_RETURN} (PS2). Computed serialized size:
\qty{1500}{\byte} for one spend and one change output, against approximately \qty{286}{\byte} for the
transparent equivalent.
\end{construction}

\paragraph{Statement proved, public inputs, witness.} Nothing beyond Sapling's own two statements
\citep{zcashprotocol}, which \fluxd{} already verifies for every shielded transaction (\shipped,
\srccode{fluxd/\allowbreak{}src/\allowbreak{}primitives/\allowbreak{}transaction.h}{155--242}, \srccode{fluxd/\allowbreak{}depends/\allowbreak{}packages/\allowbreak{}librustzcash.mk}{2--8}).

\begin{description}
  \item[Spend statement.] Public inputs $(\mathsf{rt}, \mathsf{cv}, \nf, \mathsf{rk})$; witness the
    Merkle path and position in $\tree$ (the Sapling note commitment tree, depth \num{32}) together
    with $(g_d, pk_d, v_{\note}, \mathsf{rcv}, \cm, \mathsf{rcm}, \alpha, \mathsf{ak}, \mathsf{nsk})$.
    It proves: the note commitment opens to $(g_d, pk_d, v_{\note}, \mathsf{rcm})$; $\cm$ lies in
    $\tree$ under root $\mathsf{rt}$; $\mathsf{cv}$ commits to $v_{\note}$; $\nf$ is the nullifier of
    $(\cm,\mathrm{pos})$ under $\mathsf{nk}$; $\mathsf{rk}$ re-randomises the spend authority
    $\mathsf{ak}$; and the address is diversified from $\viewkey$. Groth16 proof
    \citep{groth2016size}, \qty{192}{\byte}.
  \item[Output statement.] Public inputs $(\mathsf{cv}, \cm, \mathsf{epk})$; witness
    $(g_d, pk_d, v_{\note}, \mathsf{rcv}, \mathsf{rcm}, \mathsf{esk})$. It proves commitment and
    value-commitment integrity and $\mathsf{epk} = [\mathsf{esk}]\,g_d$. Groth16 proof,
    \qty{192}{\byte}.
  \item[Binding signature.] Under
    $\sum \mathsf{cv}_{\mathrm{spend}} - \sum \mathsf{cv}_{\mathrm{out}} -
    \mathsf{ValueCommit}(v_{\mathrm{bal}},0)$, it proves
    $\sum v_{\mathrm{spent}} - \sum v_{\mathrm{out}} = v_{\mathrm{bal}}$ exactly, with
    $v_{\mathrm{bal}}$ public.
\end{description}

\paragraph{Roles.} \emph{Prover:} the payer, or its agent. \emph{Verifiers:} every \fluxd{} node at
consensus, checking the two proofs, the binding signature and nullifier non-membership; and every
FluxOS node at the application layer, running $\mathsf{Verify}_u$ unchanged.

\paragraph{What a verifier learns.} $H_\app$ and everything $m_\app$ contains, as before; $v$ and
$v_{\mathrm{bal}}$; one nullifier per spent note, unlinkable to any commitment without $\mathsf{nk}$;
the anchor, which is a recent tree root and therefore a coarse ``this wallet had synced past height
$h$'' signal; and the change note's ciphertext. \textbf{What it does not learn:} which note was spent,
that note's value, the change value, or any key.

\begin{palgorithm}[H]
\caption{$\mathsf{Fund}_A(\principal, \app)$ --- payer side. \planned.}
\label{alg:pp-fund}
\DontPrintSemicolon
\KwIn{spending key $\mathsf{xsk}_\app$ (ZIP-32 child); registration or renewal message $m_\app$; current height $\hgt$; price table; sink address $A_{\mathrm{sink}}$}.
\KwOut{a transaction $\mathrm{tx}$ that, once confirmed, sets $e_\app \ge \hgt + \dur(m_\app)$}.
\medskip\noindent \textbf{Invariant.} $\mathrm{tx}$ has $\mathsf{vin} = \varnothing$ at all times, so Assumption~\ref{ass:pp-closed} never applies to it, and $\sum v_{\mathrm{spent}} - \sum v_{\mathrm{out}} = v_{\mathrm{bal}} = v + \mathsf{fee}$ holds by the binding signature.\medskip
$H_\app \gets \mathsf{SHA256}(m_\app)$\;
$\price \gets$ price of $m_\app$ at $\hgt$ under the height-scheduled table (\S16)\;
$v \gets \price$, or $\lceil \price/g\rceil\, g$ if a denomination grid $g$ is in force (Q6)\;
select notes under $\mathsf{xsk}_\app$ with total value $\ge v + \mathsf{fee}$\;
build \texttt{Spend\allowbreak{}Description}s; build one \texttt{Output\allowbreak{}Description} for change to a \zaddr{} derived from $\mathsf{xsk}_\app$ \tcp*{or from the next period's child (Q13)}\;
append transparent outputs $(A_{\mathrm{sink}}, v)$ and $(\mathsf{OP\_RETURN}\ H_\app, 0)$\;
set $v_{\mathrm{bal}} \gets v + \mathsf{fee}$; sign \texttt{bindingSig}\;
broadcast $\mathrm{tx}$; broadcast $m_\app$ to a FluxOS API node\;
\medskip\noindent \textbf{Failure modes.} (a) Insufficient shielded value: abort before broadcast; there is no partial-payment state. (b) $v < \price$ because the price table changed between steps 2 and 8: the payment is consumed and $e_\app$ is not advanced (the silent-underpayment defect); recovery depends on interface I4. (c) $m_\app$ does not reach a node before the funding event is scanned: the node retries the fetch up to three times, then drops the event. (d) A wallet that cannot emit an \texttt{OP\_RETURN} output alongside a Sapling spend cannot execute step 6 at all: interface I2.\medskip
\end{palgorithm}

\begin{palgorithm}[H]
\caption{$\mathsf{Verify}_u(\app, \hgt, C_{\le\hgt})$ --- node side. \shipped{} for transparent
funding events; the shielded case is covered by the same scan (interface I1, M5 resolved).}
\label{alg:pp-verify}
\DontPrintSemicolon
\KwIn{local chain prefix $C_{\le\hgt}$; local message store; the price table}.
\KwOut{$e_\app \in \mathbb{N}\cup\{\bot\}$, identical across all honest nodes with the same prefix}.
\medskip\noindent \textbf{Invariant.} No secret is read (V0); no message from another verifier is awaited (L).\medskip
\For{each confirmed $\mathrm{tx}$ at height $\hgt' \le \hgt$ with an output to $A_{\mathrm{sink}}$}{
$H \gets \mathsf{OP\_RETURN}(\mathrm{tx})$; \textbf{if} $H$ is absent \textbf{then continue}\;
fetch $m$ with $\mathsf{SHA256}(m) = H$ from local or peer message store\;
verify the owner signature on $m$; \textbf{if} invalid \textbf{then continue}\;
\textbf{if} $\mathrm{value}(\mathrm{tx}, A_{\mathrm{sink}}) \ge \price(m, \hgt')$ \textbf{then} $e_\app \gets \hgt' + \dur(m)$ \textbf{else} log a warning\;
}
\Return{} $e_\app$\;
\medskip\noindent \textbf{Failure modes.} (a) \emph{Resolved, and favourably.} The concern was that a shielded funding event would be invisible to an address-index scanner because it has an empty \texttt{vin}. It is not: \texttt{ConnectBlock}'s \texttt{vout} indexing loop runs unconditioned by \texttt{tx.vin} or coinbase status, so a z$\to$t transaction's transparent output is written to the address index regardless, and \texttt{getaddressdeltas}\slash\texttt{getaddresstxids} read it back without re-filtering \srccode{fluxd/\allowbreak{}src/\allowbreak{}main.cpp}{4072-4090}, \srccode{fluxd/\allowbreak{}src/\allowbreak{}rpc/\allowbreak{}misc.cpp}{812,1010}. A node with \texttt{-addressindex} scanning the sink address sees the payment, so Construction~\ref{con:pp-a} is not gated here --- it is gated only on reopening t$\to$z. (b) Message-store miss on all three fetch attempts: the event is dropped and the payer must re-broadcast $m_\app$. (c) A reorg past $\hgt'$: $e_\app$ is recomputed on the new prefix, which is exactly the consistency property C.\medskip
\end{palgorithm}

\paragraph{Preconditions and interface requirements.} None of these is a new cryptographic primitive.

\begin{description}
  \item[P1.] Assumption~\ref{ass:pp-closed} would have to be lifted. It will not be: the pools are
    to be retired instead (\S\ref{sec:shielded-retirement}, \Cref{rem:q1q2-resolved}), so
    Construction~\ref{con:pp-a} remains usable only by holders of pre-2021 shielded value
    (Corollary~\ref{cor:pp-who}) and only until the sunset height. \textbf{This precondition is
    unsatisfiable in the target architecture}, which is why \Cref{rem:pp-status} classes the
    constructions here as inapplicable rather than pending.
  \item[P2.] ZIP-209 turnstile enforcement, recommended on with Flux-computed checkpoints (Q2).
  \item[I1.] \texttt{processInsight} must index transactions with empty \texttt{vin} and a Sapling
    spend paying $A_{\mathrm{sink}}$ exactly as it indexes transparent ones. By construction it keys
    on \texttt{vout}; the daemon-proxy address index does return such transactions, verified in
    the indexing code \srccode{fluxd/\allowbreak{}src/\allowbreak{}main.cpp}{4072-4090} (M5, resolved).
  \item[I2.] No wallet in the corpus constructs a Sapling spend carrying an \texttt{OP\_RETURN}
    output: \SSP{} has no shielded code path, and \texttt{z\_sendmany} attaches memos only to
    \zaddr{} recipients (\shipped, \srccode{fluxd/\allowbreak{}src/\allowbreak{}wallet/\allowbreak{}rpcwallet.cpp}{3939}). Consensus
    would accept the transaction; the \texttt{Transaction\allowbreak{}Builder} needs an \texttt{OP\_RETURN} output
    path. This is small, and it is the difference between ``possible'' and ``usable''.
  \item[I3.] Under \S7's settlement nothing further is needed: PS4 divides the value of the sink
    output, which is transparent in Construction~\ref{con:pp-a}, and O7 is moot (\Cref{rem:o7-moot}).
  \item[I4.] The negative acknowledgement of \S16 is published on the FluxOS message layer keyed by
    $H_\app$. Payment disclosure (\texttt{z\_getpaymentdisclosure}, \inprogress, gated behind
    \texttt{-\allowbreak{}experimentalfeatures} plus \texttt{-\allowbreak{}paymentdisclosure},
    \srccode{fluxd/\allowbreak{}src/\allowbreak{}wallet/\allowbreak{}rpcdisclosure.cpp}{42,149}) is the right tool for the \emph{dispute}
    that follows, letting a payer prove to a chosen party what it paid. It is a disclosure tool, not a
    privacy tool, and \S20's characterisation should be read that way.
\end{description}

\subsubsection{Construction C: a shielded balance with delegated renewal}

\begin{construction}[Shielded credit balance --- \planned]
\label{con:pp-c}
The principal $\principal$ would derive a ZIP-32 child $\mathsf{xsk}_{\app}$, or
$\mathsf{xsk}_{\app,i}$ per period, and fund notes to its address --- $\zaddr \to \zaddr$ from
existing shielded value today, $\taddr \to \zaddr$ only once P1 is lifted. The ``balance'' is the
value of those notes and exists nowhere as protocol state. The renewer $\agent$ holds
$\mathsf{xsk}_{\app}$ and, at $e_\app - \Delta_{\mathrm{pre}}$, produces a
Construction~\ref{con:pp-a} transaction. The delegated budget $\budget$ is exactly the notes' value,
with scope $\scopeset = \{\text{renew }\app\text{ at no more than }\price\}$ enforced by which notes
the child key controls rather than by a protocol rule. Revocation: $\principal$ re-derives the child
and sweeps its notes, racing $\agent$ in the same way \S17 describes. Monitoring without spending:
$\principal$ retains the child's full viewing key.
\end{construction}

What holds: V0, S, C, L, D, PI against the network, RU (each renewal is an unlinkable spend; the
balance never appears on chain), and FS with per-period children. What does not hold is PI against
$\agent$, which holds the key and therefore the link. When $\agent$ is the tenant's own agent this is
not a loss and Corollary~\ref{cor:pp-agent} applies. When $\agent$ is a third-party renewer acting for
human users, it is a privacy trust point exactly as a custodial balance is, though with bounded
exposure of funds. Whether a renewer can act \emph{without} learning the link is open problem R2
(\S\ref{sec:pp-open}); this paper does not claim an answer.

\settlednote{Conditional --- inapplicable unless a shielded pool exists, and both pools are being retired (\Cref{rem:pp-status}). Q12. Renewer model for the credit path. Domain: \{own agent; third-party renewer holding
per-period child keys; custodial balance\}. Trade-off --- own agent: PI complete, but requires the
tenant to run an always-on agent; child keys: PI against the network only, with funds bounded by the
notes the child controls; custodial: neither PI nor bounded funds. Mirrors \S16's open decision on
the credit path.}

\subsection{Shape B: what a real construction would need}
\label{sec:pp-constb}

Shape B is the one place in this section where a construction outside the deployed circuits is
genuinely required, and it is stated as a \vision-maturity proposal with every parameter open. It is
given because the alternative on the table --- blind vouchers --- moves soundness off the chain, and
that trade deserves to be named rather than absorbed.

\begin{construction}[Issuer-free chain-derived entitlement --- proposed, \vision]
\label{con:pp-b}
\emph{Payment.} \texttt{OP\_RETURN} would carry \texttt{CVPN2:}~$\|~\cm$ with
$\cm = \mathsf{H}(s \| r)$ for $s, r \in \{0,1\}^{256}$ generated on-device; the amount is unchanged.
Gateways would compute \texttt{paid\_until}$[\cm]$ exactly as they compute \texttt{paid\_until}[code]
today, keyed by $\cm$ instead of by a hash of the tunnel key. Composed with
Construction~\ref{con:pp-a}, the payer is hidden on the chain as well.

\emph{Entitlement tree.} Each gateway would maintain $\tree_{\mathrm{ent}}$, a sparse Merkle tree of
depth $d_{\mathrm{ent}}$ keyed by $\cm$ with leaf
$\ell = \mathsf{H}(\cm \,\|\, \texttt{paid\_until}[\cm])$, updated in chain order. Its root
$\mathsf{rt}_\hgt$ is a deterministic function of $C_{\le\hgt}$ and the gateway's
$(\text{address},\text{price})$ configuration, hence identical on every gateway with the same chain
--- V0 and C are preserved, and L is preserved with a root-history window. The root would be published
at \texttt{/v1/info}.

\emph{Enrolment.} The client would send $(pk_t, j, \nf, \mathsf{rt}, \zkproof)$ with epoch
$j = \lfloor t / T_{\mathrm{ep}} \rfloor$ and $\nf = \mathsf{PRF}_s(j)$, where $\zkproof$ is a proof
for the relation
\begin{multline*}
  \mathcal{R}_{\mathrm{ent}} = \Big\{\,
    \big((\mathsf{rt}, j, \nf, \eta);\ (s, r, w, \mathrm{path})\big) \ :\
    \mathsf{Leaf}(\mathsf{H}(s\|r), w) \in \tree_{\mathrm{ent}} \text{ under } \mathsf{rt} \\
    {}\ \wedge\ w \ge j\,T_{\mathrm{ep}}
    \ \wedge\ \nf = \mathsf{PRF}_s(j)
  \,\Big\}
\end{multline*}
with public inputs $(\mathsf{rt}, j, \nf, \eta)$ where $\eta = \mathsf{H}(pk_t)$, witness
$s, r \in \{0,1\}^{256}$, $w \in \mathbb{N}$ (the paid-through time), and
$\mathrm{path} \in (\{0,1\}^{256})^{d_{\mathrm{ent}}}$. The public input $\eta$ is bound into the
proof to prevent proof theft --- the Semaphore ``signal'' pattern.

\emph{Gateway.} Accept if $\mathsf{rt}$ is among the gateway's last $N_{\mathsf{rt}}$ published roots
and the proof verifies; set \texttt{tier}$[pk_t] \gets$ premium; key the rate limiter by $\nf$, so
that all tunnel keys presenting the same $\nf$ within epoch $j$ share one limiter --- \emph{the same
sharing bound the per-\texttt{pubkey} limiter enforces today}
(\shipped, \srccode{cumulusvpn/\allowbreak{}gateway/\allowbreak{}internal/\allowbreak{}limiter/\allowbreak{}limiter.go}{105--123}). D therefore needs no
global spent set: entitlement is time-based rather than consumable, and sharing across gateways is
already permitted in the shipped model.

\emph{Epoch rotation.} A fresh $pk_t$ and a fresh $\nf$ per epoch make sessions unlinkable across
epochs and linkable within one.
\end{construction}

\paragraph{What the gateway learns.} $(pk_t, \nf, j, \mathrm{IP}, \text{traffic})$ and the fact that
some active subscriber is present. Not $\cm$, not the payment, not the payer; linking $\nf$ to $\cm$
requires $s$. The anonymity set is the set of leaves with $w \ge j\,T_{\mathrm{ep}}$ --- every active
subscriber --- refined downward by first-use timing correlation between a new leaf appearing on chain
and a new enrolment arriving (\S\ref{sec:pp-attacks}), whose magnitude is CumulusVPN's payment rate
and is unmeasured (M3).

\paragraph{What it requires that Flux does not have.} Four things, and the interfaces are stated so
the construction can be evaluated even though it is not built.

\begin{enumerate}
  \item \textbf{The circuit.} Approximately $d_{\mathrm{ent}}$ hash evaluations for the Merkle path,
    two for the leaf, one PRF evaluation and one range check --- of order $10^4$ constraints with a
    SNARK-friendly hash, which is small next to Sapling's spend circuit.
  \item \textbf{A proving key.} Under Groth16 this means a circuit-specific ceremony
    \citep{groth2016size,bowe2017mpc}, which would be the first ceremony this project has ever run:
    \S20 documents that \fluxd{} re-hosts the 2018 Zcash Sapling parameters verbatim and performs no
    ceremony of its own (\shipped, \srccode{fluxd/\allowbreak{}zcutil/\allowbreak{}fetch-params.sh}{220--222},
    \srccode{fluxd/\allowbreak{}src/\allowbreak{}init.cpp}{754--806}).
  \item \textbf{Or a move to a ceremony-free system.} Halo-class recursive arguments remove the
    trusted setup \citep{bowe2019halo}. \textbf{Two properties must be kept strictly apart here:}
    no-trusted-setup is not post-quantum. Halo2 rests on discrete-log hardness on a pairing-friendly
    or Pasta-style curve and is no more post-quantum than Groth16 over BLS12-381. Removing the
    inherited ceremony serves the decentralization requirement of \S20; it does not serve the
    post-quantum requirement of \S8, which by settled scope excludes the shielded layer entirely.
  \item \textbf{Client and gateway code.} A WASM prover in the TypeScript and React Native clients,
    and a BLS12-381 verifier in the Go gateway.
\end{enumerate}

\unsupported{Proving time for $\mathcal{R}_{\mathrm{ent}}$ in a browser or mobile client is asserted
nowhere measurable; the ``sub-second'' figure in the mechanism design is an order-of-magnitude
inference from circuit size, not a measurement. Either measure it on target client hardware or drop
the claim.}

\settlednote{Conditional --- inapplicable unless a shielded pool exists, and both pools are being retired (\Cref{rem:pp-status}). Q8. Proving system for any new Flux circuit. Domain: \{Groth16 with a Flux-run ceremony;
a Halo2/PLONK-class system with universal or no setup, undated; defer shape B\}. Trade-off: every new
Groth16 statement needs its own ceremony and this project has never run one; a ceremony-free system
removes that political cost and adds a proving-system migration that is scheduled nowhere.}

\settlednote{Conditional --- inapplicable unless a shielded pool exists, and both pools are being retired (\Cref{rem:pp-status}). Q10 and Q11, construction B parameters. Epoch length $T_{\mathrm{ep}} \in [\qty{1}{\hour},
\qty{30}{\day}]$: shorter gives less intra-epoch linkability at the cost of more re-enrolment and
proving; longer lets the nullifier $\nf$ decay into a stable pseudonym. Root-history window
$N_{\mathsf{rt}} \in [10, 10^4]$ roots: larger tolerates staler clients at the cost of gateway state.
Tree depth $d_{\mathrm{ent}}$ is fixed by the target subscriber population, which is unmeasured (M3).}

\begin{remark}[The voucher alternative, and what it costs]
\label{rem:pp-voucher}
CumulusVPN's stated v2 is blind-signature vouchers with quorum signing (\vision,
\srcdoc{cumulusvpn/docs/04-payments.md:84--87}). Blind signatures give unlinkability by construction
and, for the Chaum--RSA variant, information-theoretically. But they require an \emph{issuer} holding
a secret key, and nothing on the chain binds tokens issued to payments received: \textbf{soundness
moves from the chain to the issuer.} Definition~\ref{def:pp-obligation}'s V0 forbids a secret in
\emph{verification} only, so an issuer is admissible --- but it becomes the party whose honesty S
depends on. Threshold or quorum issuance bounds the availability risk and does not restore
chain-anchored soundness. That is the real trade, and both options should be presented with it:
Construction~\ref{con:pp-b} keeps soundness on the chain and pays for it with a circuit and a
ceremony; vouchers avoid the circuit and pay for it with an issuer.
\end{remark}

\subsection{The anonymity set, measured}
\label{sec:pp-anonymity}

No construction manufactures an anonymity set. This subsection measures the one that exists, and the
numbers are the section's most load-bearing empirical content.

\subsubsection{The population is not the payer's anonymity set}

For shape A the payment names $H_\app$, so the \emph{workload-side} anonymity set is $1$ by
construction and the \num{1195} deployed applications are not the anonymity set of anything. The
question that matters is whether the (amount, time) pair leaks more than the name already does. For
public-content applications it does not, because $\price$ is a deterministic function of the public
specification. Table~\ref{tab:pp-classes} gives the partition structure.

\begin{table}[H]
\centering
\small
\caption{Price-class structure over the \num{1195} deployed applications
\srcmeasured{api.runonflux.io/\allowbreak{}apps/\allowbreak{}globalappsspecifications}{2026-09-03}. The tuple partition
$(\mathrm{cpu},\mathrm{ram},\mathrm{hdd},\ninst,\dur,\mathsf{enterprise},\mathsf{staticip},
\mathsf{nodes})$ is a lower bound on class size under \emph{any} deterministic price function; the
period-price partitions are upper bounds up to rare collisions.}
\label{tab:pp-classes}
\begin{tabular}{@{}lrrrr@{}}
\toprule
\textbf{Partition} & \textbf{Classes} & \textbf{Singletons} & \textbf{Median class} & \textbf{Largest} \\
\midrule
resource tuple (lower bound)        & 778 & 698 (\qty{58.4}{\percent}) & 1 & 57 \\
period price, base only (lower bound) & 349 & 201 (\qty{16.8}{\percent}) & 7 & 63 \\
period price, $\ninst/3$ variant (upper bound) & 381 & 229 (\qty{19.2}{\percent}) & 6 & 62 \\
\textbf{period price, \Cref{def:price} as corrected} & \textbf{370} & \textbf{222} (\qty{18.6}{\percent}) & \textbf{8} & \textbf{61} \\
monthly price only, ignoring $\dur$ & 170 & \phantom{0}86 (\qty{7.2}{\percent})  & 66 & 210 \\
\addlinespace
price \emph{and} same block         & 1{,}195 & 1{,}195 (\qty{100}{\percent}) & 1 & 1 \\
price and same hour                 & --- & 976 (\qty{81.7}{\percent}) & 1 & 11 \\
price and same day                  & --- & 669 (\qty{56.0}{\percent}) & 1 & 18 \\
price and same week                 & --- & 439 (\qty{36.7}{\percent}) & 2 & 35 \\
\bottomrule
\end{tabular}
\end{table}

If every application renews on schedule, renewal volume is \num{0.0123} per block, \num{1.48} per
hour, \textbf{\num{35.4} per day} and \num{248} per week; the next 30 days carry a median of \num{24}
and a maximum of \num{98} expiries per day. That figure is an upper bound, because churn is unmeasured
(M2). The median period price is \num{1.13}\flux, the range \numrange{0.02}{79.85}\flux;
\num{459} applications (\qty{38.4}{\percent}) sit at the default \num{88000}-block period and
\num{107} (\qty{9.0}{\percent}) at the \num{1056000}-block maximum.

The reading is this. \textbf{At \num{35.4} renewals per day a renewal is a singleton within its block
\qty{100}{\percent} of the time and within its day \qty{56.0}{\percent} of the time.} A payment ``of a
distinctive size just before a distinctive deployment'' is not an attack scenario, it is the normal
case. For shape A that identifies the application, which the \texttt{OP\_RETURN} already does, and
says nothing about the payer. For any coupon-style or Construction~\ref{con:pp-b} scheme paying exact
amounts, it identifies the entitlement and defeats the scheme --- which is why such schemes require
\emph{denominated} amounts, at a measurable cost.

\begin{table}[H]
\centering
\small
\caption{Denomination trade-off, computed on Flux's own application population. Overpayment is
expressed as a share of list revenue; under \PNR{} it would flow to the Foundation and the fee pool
by the standard split, i.e.\ a privacy fee paid to the network.}
\label{tab:pp-denom}
\begin{tabular}{@{}lrrrrr@{}}
\toprule
\textbf{Grid $g$} & \textbf{Classes} & \textbf{Singletons} & \textbf{Median class} & \textbf{Largest} & \textbf{Overpayment} \\
\midrule
none (exact)      & 349 & 201 (\qty{16.8}{\percent}) & 7 & 63 & \qty{0}{\percent} \\
\num{0.5}\flux    & 43 & 15 (\qty{1.3}{\percent}) & 118 & 291 & \qty{9.0}{\percent} \\
\num{1}\flux      & 30 & 13 (\qty{1.1}{\percent}) & 227 & 559 & \qty{17.3}{\percent} \\
\num{2}\flux      & 21 & \phantom{0}7 (\qty{0.6}{\percent}) & 786 & 786 & \qty{39.0}{\percent} \\
powers of two     & 13 & \phantom{0}0 & 227 & 268 & \qty{51.4}{\percent} \\
\num{5}\flux      & 12 & \phantom{0}4 (\qty{0.3}{\percent}) & 1{,}037 & 1{,}037 & \qty{119.7}{\percent} \\
\bottomrule
\end{tabular}
\end{table}

\begin{figure}[H]
\centering
\begin{tikzpicture}
\begin{axis}[
  width=0.80\linewidth, height=6.2cm,
  xlabel={denomination grid $g$ (FLUX); $g=0$ denotes exact pricing},
  ylabel={median price-class size},
  ymode=log, log basis y={10},
  xmin=-0.3, xmax=5.3, ymin=4, ymax=3000,
  xtick={0,0.5,1,2,5},
  xticklabels={exact,0.5,1,2,5},
  axis y line*=left, axis x line*=bottom,
  grid=major, major grid style={draw=gray!20},
  mark size=2pt,
]
\addplot[thick, mark=*] table[x index=0, y index=1] {../data/private_renewal_denom.dat};
\end{axis}
\begin{axis}[
  width=0.80\linewidth, height=6.2cm,
  axis y line*=right, axis x line=none,
  ylabel={overpayment (\% of list revenue)},
  xmin=-0.3, xmax=5.3, ymin=0, ymax=130,
  mark size=2pt,
]
\addplot[thick, dashed, mark=square] table[x index=0, y index=2] {../data/private_renewal_denom.dat};
\end{axis}
\end{tikzpicture}
\caption{The denomination trade-off measured on the \num{1195}-application population. Solid, left
axis, logarithmic: median price-class size. Dashed, right axis, linear: overpayment as a share of
list revenue. A \num{1}\flux{} grid buys a median class of \num{227} at \qty{17.3}{\percent}
overpayment; a \num{0.5}\flux{} grid buys \num{118} at \qty{9.0}{\percent}. This is the classical
fixed-denomination trade-off, evaluated on Flux's own demand rather than assumed.}
\label{fig:pp-denom}
\end{figure}

\settlednote{Conditional --- inapplicable unless a shielded pool exists, and both pools are being retired (\Cref{rem:pp-status}). Q6. Denomination grid $g \in \{\text{none}, 0.5, 1, 2, 5\}$\flux{} or powers of two, plus
the destination of the overpayment (\PNR{} split, or credited forward to the next period). Trade-off
is Table~\ref{tab:pp-denom} exactly: median class $7 \to 118 \to 227 \to 786$ against overpayment
$0 \to \qty{9.0}{\percent} \to \qty{17.3}{\percent} \to \qty{39.0}{\percent}$. Note that denomination
buys nothing for shape A, where the workload is named anyway; it is a shape-B and coupon-scheme
parameter.}

\subsubsection{The payer's anonymity set is a closed pool of unknown membership}

\textbf{This is the number the section must print.} The payer-side anonymity set for
Construction~\ref{con:pp-a} is not the application population; it is the holder set of the Sapling
pool: \num{77188.75921324}\flux{} of value (plus \num{19291.17735041}\flux{} in Sprout), closed since
height \num{835554}, with \emph{note count and holder count unmeasured}. The
\texttt{commitments} field of \texttt{getblockchaininfo} (\num{815736} at tip) is, in the zcashd-2.x
lineage, the \emph{Sprout} tree size and does not answer the question.

\begin{openproblem}[M1: the effective size of a closed pool]
\label{op:pp-m1}
Determine the number of unspent Sapling notes at the tip --- tree size minus nullifier-set size ---
and any defensible holder-count estimate that chain analysis permits. Until this is measured, the
paper claims the \emph{mechanism} of Construction~\ref{con:pp-a} and does not claim realised privacy
for it. Interpreting the measurement (which notes are plausibly spendable, how holders cluster) is
chain-analysis research with no standard answer, and is listed separately as R3.
\end{openproblem}

\subsection{Cost, and the binding limit on shielded throughput}
\label{sec:pp-cost}

\begin{table}[H]
\centering
\small
\caption{Serialized sizes, derived from the deployed constants (\shipped,
\srccode{fluxd/\allowbreak{}src/\allowbreak{}flux/\allowbreak{}Zelcash.h}{8--24}, \srccode{fluxd/\allowbreak{}src/\allowbreak{}flux/\allowbreak{}JoinSplit.hpp}{22--26},
\srccode{fluxd/\allowbreak{}src/\allowbreak{}primitives/\allowbreak{}transaction.h}{155--242}).}
\label{tab:pp-sizes}
\begin{tabular}{@{}lrl@{}}
\toprule
\textbf{Object} & \textbf{Bytes} & \textbf{Derivation} \\
\midrule
Sapling note plaintext / \texttt{encCiphertext} & 564 / 580 & $1+11+8+32+512$, then $+16$ \\
\texttt{outCiphertext} & 80 & $32+32+16$ \\
\texttt{Spend\allowbreak{}Description} & 384 & $4\times 32 + 192 + 64$ \\
\texttt{Output\allowbreak{}Description} & 948 & $3\times 32 + 580 + 80 + 192$ \\
Construction~\ref{con:pp-a} renewal & \textbf{1{,}500} & header, counts, two vouts, one spend, one output, \texttt{bindingSig} \\
transparent renewal (1 P2PKH input) & $\approx 286$ & --- \\
Construction~\ref{con:pp-b} proof (off-chain) & 192 & Groth16 \\
\bottomrule
\end{tabular}
\end{table}

A shielded renewal is \qty{1500}{\byte} against approximately \qty{286}{\byte} transparent, a factor
of \num{5.24}. With \texttt{MAX\_BLOCK\_SIZE} $=$ \qty{2000000}{\byte}, block capacity is
$\lfloor 2{,}000{,}000/1{,}500 \rfloor = \num{1333}$ such transactions per block, \num{44.4} per
second, implying \num{2666} Groth16 verifications per block.

\begin{property}[Block space is not the constraint]
\label{prop:pp-space}
At measured demand of \num{0.0123} renewals per block, Construction~\ref{con:pp-a} would consume
\qty{0.0009}{\percent} of \texttt{MAX\_BLOCK\_SIZE}. At \numlist{10;100;1000} times measured demand
the figures are \qty{0.01}{\percent}, \qty{0.09}{\percent} and \qty{0.92}{\percent}. Even if all
\num{7486} requested instances were re-registered as separate applications \emph{daily}, the load
would be \num{2.6} transactions per block.
\end{property}

\begin{property}[Verification is the constraint]
\label{prop:pp-verify}
Let $t_{G16}$ be single-core Groth16 verification time and $N_{\mathrm{proofs}}$ the proof count in a
block; verification cost is $T_{\mathrm{verify}} = N_{\mathrm{proofs}}\cdot t_{G16}$ on one core. A
block filled with Construction~\ref{con:pp-a} transactions carries $N_{\mathrm{proofs}} = \num{2666}$;
a pathological all-\texttt{Spend\allowbreak{}Description} block carries
$\lfloor 2{,}000{,}000/384\rfloor = \num{5208}$. Against the \qty{30}{\second} \PoN{} interval:
\[
  \num{2666}\cdot t_{G16} > \qty{30}{\second} \iff t_{G16} > \qty{11.3}{\milli\second},
  \qquad
  \num{5208}\cdot t_{G16} > \qty{30}{\second} \iff t_{G16} > \qty{5.8}{\milli\second}.
\]
Concretely, at \qtylist{2;5;10}{\milli\second} the pathological block takes
\qtylist{10.4;26.0;52.1}{\second} to verify. \textbf{\texttt{MAX\_BLOCK\_SIZE} was sized for
transparent transactions; for shielded traffic the binding limit on throughput is single-core proof
verification, not block space.} At ordinary demand the point is moot --- at $1000\times$ demand,
\num{24.6} proofs per block, or \qty{0.25}{\second} at \qty{10}{\milli\second} --- but the
adversarially-full block would become reachable once shielded traffic was encouraged, which is what a
reopening of $\taddr \to \zaddr$ would have done. With reopening cancelled (\S\ref{sec:pp-blocker}),
no such activation height arrives and the bound remains a ceiling on a pool that can only drain.
\end{property}

\unsupported{$t_{G16}$ on Cumulus, Nimbus and Stratus hardware is a parameter here, not a
measurement (M4). It must be measured with \texttt{zcbenchmark}, whose \texttt{verifysaplingspend}
and \texttt{verifysaplingoutput} types the fork retains
(\srccode{fluxd/\allowbreak{}src/\allowbreak{}wallet/\allowbreak{}rpcwallet.cpp}{2972--2979}). Likewise, the claim that the 2018-era \texttt{librustzcash} pin
\texttt{06da3b9a} predates batch verification is an inference from the pin date, not from reading the
crate.}

\settlednote{Conditional --- inapplicable unless a shielded pool exists, and both pools are being retired (\Cref{rem:pp-status}). Q17. Response to Property~\ref{prop:pp-verify}. Domain: \{no change; a consensus cap
$N_{\mathrm{proofs}}$ per block; a \texttt{librustzcash} upgrade enabling batch verification\}.
Trade-off: a cap is a one-constant consensus change that bounds worst-case validation time and
also caps legitimate shielded throughput; an upgrade removes the cap's need but moves the fork's
pinned cryptography forward by several years of upstream changes, which is a much larger review
surface. Only a reopening of $\taddr \to \zaddr$ would have let shielded volume grow past today's
zero, and reopening is cancelled (\S\ref{sec:pp-blocker}).}

Two further cost notes. Verifying keys are already present on every node --- \texttt{ZC\_LoadParams}
refuses to start without all three parameter files (\shipped,
\srccode{fluxd/\allowbreak{}src/\allowbreak{}init.cpp}{754--806}) --- so \num{6489} nodes verify Sapling proofs today at zero
marginal deployment cost. And Sprout value, the \num{19291.18}\flux{} half of the pool, requires the
\qty{725}{\mega\byte} \texttt{sprout-groth16.\allowbreak{}params} and gigabyte-class proving memory; the paper
treats it as legacy.

\subsection{Invariants}
\label{sec:pp-invariants}

Table~\ref{tab:pp-invariants} states which invariants hold under which construction. Each is stated so
that a violation is observable.

\begin{table}[H]
\centering
\small
\caption{Invariants by construction. A: shielded direct payment; C$_{\mathrm{own}}$: credit balance
renewed by the payer's own agent; C$_{3\mathrm{p}}$: credit balance renewed by a third party;
B: issuer-free entitlement.}
\label{tab:pp-invariants}
\begin{tabular}{@{}llcccc@{}}
\toprule
& \textbf{Invariant} & \textbf{A} & \textbf{C}$_{\mathrm{own}}$ & \textbf{C}$_{3\mathrm{p}}$ & \textbf{B} \\
\midrule
I1 & Conservation: $\sum v_{\mathrm{spent}} - \sum v_{\mathrm{out}} = v_{\mathrm{bal}}$, pool value falls by exactly $v_{\mathrm{bal}}$ & yes & yes & yes & n/a \\
I2 & Non-negativity of $\pool$, $e_\app$ and note values & yes & yes & yes & yes \\
I3 & No double-funding from one payment (D) & yes & yes & yes & vacuous \\
I4 & No forgery of funded status (S) & yes$^\dagger$ & yes$^\dagger$ & yes$^\dagger$ & yes \\
I5 & No linkage across renewals of one deployment (RU) & yes$^\ddagger$ & yes$^\ddagger$ & \emph{no} & yes \\
I6 & Forward secrecy of the payer across periods (FS) & partial & partial & \emph{no} & yes \\
I7 & Liveness under partition (L) & yes & yes & yes & yes$^\S$ \\
I8 & No double-claim across an instance migration & yes & yes & yes & n/a \\
\bottomrule
\end{tabular}
\end{table}

\begin{description}
  \item[I1, conservation.] Construction~\ref{con:pp-a} moves value through the existing Sapling
    turnstile arithmetic: the binding signature forces
    $\sum v_{\mathrm{spent}} - \sum v_{\mathrm{out}} = v_{\mathrm{bal}}$, and the pool's
    \texttt{chainValue} falls by exactly $v_{\mathrm{bal}}$. The coinbase is untouched, so \S7's
    bounded-coinbase requirement is trivially preserved: this mechanism adds no coinbase output.
  \item[I2, non-negativity.] A $\zaddr \to \taddr$ transaction with $v_{\mathrm{bal}}$ exceeding the
    pool's value is rejected \emph{only when turnstile enforcement is on}; today it is monitored and
    not enforced (Q2). The invariant is conditional on that decision and the paper states the
    dependency rather than the invariant alone.
  \item[I3, no double-funding.] $H_\app$ is resolved once --- \texttt{check\allowbreak{}And\allowbreak{}Request\allowbreak{}App} returns
    early when the permanent message already exists --- so a second payment carrying an
    already-resolved hash is \emph{ignored}, that is, lost. This joins the silent-underpayment defect
    as a failure mode that interface I4 must report. Under Construction~\ref{con:pp-c} each renewal
    is a distinct $m_\app$ with a distinct hash. Under Construction~\ref{con:pp-b} the nullifier keys
    the limiter and D is vacuous, because entitlement is time-based rather than consumable.
  \item[I4, no forgery of funded status.] Requires a confirmed transfer bound to $H_\app$; forging one
    requires forging a Sapling proof or the binding signature. $\dagger$: under Groth16 with the
    \emph{inherited 2018 parameters}, this is exactly the counterfeiting scenario \S20 documents, and
    it is bounded to the pool's declared value only if turnstile enforcement is on.
  \item[I5, no linkage across renewals.] Consecutive Construction~\ref{con:pp-a} transactions share
    the name (via each period's $H_\app$), the amount (same price) and the periodicity --- all public
    by design --- and nothing else: fresh $\nf$, $\mathsf{cv}$, $\mathsf{rk}$ and ciphertext.
    $\ddagger$ marks the residual: if period $i$'s change note is spent in period $i+1$ and the
    period-$(i{+}1)$ child key is compromised, the adversary decrypts that note and learns it was
    created by period $i$'s renewal --- linking two renewals that $H_\app$ already links, and
    revealing how many periods the sub-wallet has funded. No payer identity is exposed. Forwarding
    change to the next period's child confines the residual to one step (Q13).
  \item[I6, forward secrecy.] With per-period ZIP-32 children and parents discarded on the agent,
    compromise at period $i$ yields keys for periods $\ge i$ only. ``Partial'' because the chain
    retains every ciphertext forever, so compromise of the principal's \emph{master} seed reveals all
    periods --- inherent to hierarchical-deterministic shielded wallets and shared with Zcash.
  \item[I7, liveness under partition.] $\mathsf{Verify}_u$ depends only on the local prefix; a
    partitioned node evaluates on its own view and reconciles by the same disconnect logic reorgs use
    today. $\S$: for Construction~\ref{con:pp-b}, a gateway with a stale root rejects proofs against
    newer roots and the client re-proves against the published root, one extra round trip, with
    $N_{\mathsf{rt}}$ bounding staleness tolerance.
  \item[I8, no double-claim across migration.] Funded status is chain-derived and node-independent, so
    when FluxOS redeploys instances to other nodes there is no per-node claim to duplicate. Across the
    dead Sprout$\to$Sapling migration path, Sprout value remains spendable by a joinsplit with
    \texttt{vpub\_new} and empty \texttt{vin} --- therefore not caught by
    Assumption~\ref{ass:pp-closed} --- at higher proving cost; nothing is claimed twice.
\end{description}

\subsection{Attack analysis}
\label{sec:pp-attacks}

\paragraph{Timing and amount correlation.} Renewal heights are public and periodic,
$\hgt_0 + i\,\dur$. An adversary with a network vantage watches for a shielded transaction entering
the mempool in the window before each renewal height of $\app$ and records the first-seen peer. Over
$i$ renewals the intersection of first-seen peer sets converges on the payer's broadcast point.
The mitigation is behavioural, not cryptographic: prepay, which the deployed
\texttt{max\allowbreak{}Blocks\allowbreak{}Allowance} of \num{1056000} blocks (\num{366.7} days) already permits and which
reduces events from \num{11.95} to \num{1.00} per year; randomise the submission lead within
$[\Delta_{\mathrm{pre}}, \dur]$; and broadcast through a node the tenant does not otherwise use --- or
through the VPN the network itself sells. \fluxd{} has no Dandelion-style relay and inherits
Bitcoin-era transaction propagation, so this attack is not addressed at the P2P layer.

\paragraph{The intersection attack across renewals.} On chain, Construction~\ref{con:pp-a} gives the
adversary nothing to intersect: each renewal exposes a fresh nullifier and no persistent identifier.
The attack lives entirely at the network layer, as above. Against a payer broadcasting from a fixed
vantage it converges in a handful of periods; against one who prepays a year it has a single sample.
The interaction with \S23 should be stated: a tenant who prepays a year privately --- the cheapest
mitigation --- carries a larger forfeitable remainder if the application is later removed by the
moderation quorum. At the median period price that is \num{13.5}\flux{} for a year; at the maximum
measured period price of \num{79.85}\flux{} it is \num{955}\flux.

\paragraph{A node operator who is also the payer.} On chain, no advantage: its own transactions look
like anyone's. At the FluxOS layer it can submit its own registrations through its own node and
broadcast its own transactions from it, eliminating the metadata trust point \emph{for itself};
conversely, as an adversary it observes the metadata of every tenant that uses its nodes, and with
\num{237} nodes under one identity that is \qty{3.7}{\percent} of API entry points. Self-dealing gains
nothing economically --- \S7's timing recapture is structural and independent of whether the fee's
origin is visible --- but privacy removes the \emph{audit trail} that would let anyone notice a
wash-deployment pattern. The defence against wash deployment is \PNR's economics, not observability,
and \S7 should be read as carrying that weight.

\paragraph{What the transparent layer leaks even when the payment is shielded.} Measured on the
specification set itself, which no payment scheme touches
\srcmeasured{api.runonflux.io/\allowbreak{}apps/\allowbreak{}globalappsspecifications}{2026-09-03}:

\begin{itemize}
  \item \textbf{Owner pseudonyms.} \num{659} distinct \texttt{owner} identifiers; \num{588} own
    exactly one application; the largest owns \num{246} (\qty{20.6}{\percent}) and is a
    Foundation-operated address. Excluding it, \num{949} applications across \num{658} owners, of whom
    \num{15} hold five or more (\num{223} applications in total). Owner-key reuse links applications
    across renewals and across deployments \emph{regardless of payment privacy}. The agent-native
    default should be one owner key per application (Q15).
  \item \textbf{Contacts.} \num{688} applications (\qty{57.6}{\percent}) publish at least one
    non-blank contact string in the public specification. Any such string is a direct identity leak
    that dominates payment privacy for the application carrying it (Q14).
  \item \textbf{The registration message} carries the owner's signature and reaches the network
    through one API node, which sees the client's IP address and is identifiable to peers as the
    broadcast origin. This is a metadata trust point independent of the chain (Q7).
  \item \textbf{Amount and timing}, per \S\ref{sec:pp-anonymity}; the \textbf{anchor}, a coarse wallet
    sync height; the \textbf{fee}; and \textbf{\texttt{valueBalance}}, which equals $v + \mathsf{fee}$
    and is public anyway.
\end{itemize}

\paragraph{Griefing and denial of service.} Invalid Sapling proofs cost the submitter nothing --- they
are rejected before any fee --- and cost each verifying node one Groth16 verification; this is bounded
by P2P DoS scoring as in Zcash. Block-filling is Property~\ref{prop:pp-verify}. For
Construction~\ref{con:pp-b}, a gateway that publishes a root it will not accept denies service to
anonymous clients only, leaves free-tier clients unaffected, and is detectable by clients but not by
the network.

\paragraph{Byzantine and colluding node sets.} A byzantine minority cannot forge funded status
(I4: every node verifies independently) and cannot deny it (L); it can withhold or delay a tenant's
registration message, as it can today, with the client retrying against another node; it can log
metadata at the nodes it runs. A colluding majority can rewrite the chain, at which point the
funded-status question is moot. Short of that, a majority of \emph{API entry points} still cannot
break PI cryptographically; it widens the metadata view above. A colluding majority of
Construction~\ref{con:pp-b} gateways learns nothing more than one gateway does, because gateways hold
no secret --- which is the point of making the construction issuer-free.

\paragraph{Sustained low demand.} This is the operative failure mode and it is present now:
\num{35.4} renewals per day, of which zero are shielded. Privacy degrades continuously with volume;
there is no threshold below which it stops working, only a shrinking $\mathsf{AS}$. The paper does not
claim privacy at current volume. It claims the mechanism and prints the population figures.

\paragraph{Post-quantum scope.} Every construction here rests on BLS12-381 pairings and Jubjub
discrete logarithms, and none is post-quantum --- by the settled scope decision of \S8, which places
the shielded layer outside post-quantum migration. The one interaction worth stating: a quantum
adversary breaks \emph{hiding retroactively}, since it can decrypt stored ciphertexts and recover
keys, so PI for payments made today does not survive that adversary. This is Zcash's position as well,
and \S20 states it.

\subsection{Open problems, and open parameters}
\label{sec:pp-open}

The two categories are kept apart deliberately. A research problem has no known good answer, or no
answer without a new trust assumption; an engineering item is a substantial build with no unsolved
question. \S\ref{sec:feasibility} must class them separately and must not collapse them.

\begin{openproblem}[R1: unlinkable exact-amount payment at low volume]
\label{op:pp-r1}
With \num{35.4} renewals per day and \num{349} price classes, the pair (amount, time) is identifying
(Table~\ref{tab:pp-classes}). The only known remedies are denomination, which costs overpayment
(Table~\ref{tab:pp-denom}), and batching or delay, which costs latency; both are policies, not
constructions. Whether a mechanism exists that is simultaneously exact-amount, unlinkable and
\PNR-compatible is open. Property~\ref{prop:pp-amount} shows that amount hiding is not that
mechanism. \textbf{This is a property of the population, not of any construction} --- which is why no
circuit fixes it.
\end{openproblem}

\begin{openproblem}[R2: third-party renewal without the renewer learning the link]
\label{op:pp-r2}
Construction~\ref{con:pp-c} hands the renewer the spending key and therefore the link. A renewer that
spends value bound to $\app$ without learning which principal funded it is the
delegatable-anonymous-credential / anonymous-e-cash-with-delegation setting. No issuer-free instance
for \emph{this} setting --- many independent verifiers, no shared secret, soundness anchored on the
chain --- is known to this author, and this paper does not invent one. For autonomous agents the
problem is moot (Corollary~\ref{cor:pp-agent}); it binds only for human users delegating to a
third-party renewal service.
\end{openproblem}

\begin{openproblem}[R3: the effective anonymity set of a closed pool]
\label{op:pp-r3}
M1 (Open Problem~\ref{op:pp-m1}) is a measurement. \emph{Interpreting} it --- which notes are
plausibly still spendable, how holders cluster, what a realistic adversary can exclude --- is
chain-analysis research with no standard answer, and it is made harder, not easier, by the pool being
closed: a pool that can only lose value admits elimination arguments that a growing pool does not.
\end{openproblem}

Everything else is engineering: (E1) the P1 network upgrade and P2 turnstile activation with
Flux-computed checkpoints; (E2) interfaces I1 and I3, the FluxOS scanner and \PNR{} validation for
empty-\texttt{vin} payments; (E3) interface I2, a \texttt{Transaction\allowbreak{}Builder} \texttt{OP\_RETURN} path
and a wallet or agent SDK that uses it; (E4) interface I4, the negative acknowledgement on the message
layer with payment disclosure as the dispute tool; (E5) Construction~\ref{con:pp-c} as a ZIP-32
child-key delegation object composed with \S17's delegation primitive; (E6) Construction~\ref{con:pp-b}
--- circuit, ceremony or Halo2, client prover, gateway verifier; and (E7) the denominated-pricing
option and the response to Property~\ref{prop:pp-verify}. \textbf{The feasibility assessment of \S\ref{sec:feasibility}
should therefore read: shape A is Engineering; shape C with the payer's own agent is Engineering;
shape B is Engineering plus a ceremony; and the Research content is R1--R3, of which R1 is a property
of the population rather than of any construction.}

\begin{table}[H]
\centering
\scriptsize
\caption{Decisions, open unless marked. None is guessed anywhere in this section; each is mirrored in the
decisions register. ``Affects'' names the construction.}
\label{tab:pp-open}
\begin{tabular}{@{}llp{0.26\linewidth}p{0.36\linewidth}@{}}
\toprule
\textbf{Id} & \textbf{Affects} & \textbf{Domain} & \textbf{Trade-off} \\
\midrule
Q1 & A, C & withdrawn: no reopening (\Cref{rem:q1q2-resolved}) & pool unreachable until activation versus turnstile-bound soundness at reopening \\
Q2 & A, C & answered: on, for the drain window (\Cref{rem:q1q2-resolved}) & bounded soundness failure versus none \\
Q3 & A & settled: $H(m_\app)$ in \texttt{OP\_RETURN} \srcintent{08-answers-round-2.md, round 2} & binding $H(m_\app)$ is what nodes must act on; blinding belongs to shape B \\
Q4 & A & moot by retirement: transparent only (\S7, \Cref{rem:o7-moot}) & exclusion makes this section vacuous; hidden amount is impossible (Prop.~\ref{prop:pp-amount}) \\
Q5 & A & amount hiding for non-\PNR{} classes: none / an exempt class & an exempt class restores AH for enterprise workloads at the cost of revenue outside \S7's split \\
Q6 & A, B & denomination grid, Table~\ref{tab:pp-denom} & median class versus overpayment \\
Q7 & A, C & submission path: any node / own node / VPN / no-logging rule & metadata trust point; a logging rule is policy and unenforceable \\
Q8 & B & Groth16 with a ceremony / Halo2-class / defer & first-ever ceremony versus an undated proving-system move \\
Q9 & B & Construction~\ref{con:pp-b} / single issuer / threshold issuer & soundness on chain versus soundness at an issuer \\
Q10 & B & $T_{\mathrm{ep}} \in [\qty{1}{\hour}, \qty{30}{\day}]$ & intra-epoch linkability versus re-enrolment cost \\
Q11 & B & $N_{\mathsf{rt}} \in [10, 10^4]$ & stale-client liveness versus gateway state \\
Q12 & C & own agent / third-party renewer / custodial & PI complete / PI versus the network only / neither \\
Q13 & C & change forwarding to per-period children: yes / no & free on chain; confines I5's residual to one step \\
Q14 & all & contacts: keep / encrypt / strip from the hashed object & \num{688} applications leak an identifier today and no payment scheme fixes it \\
Q15 & A & one owner key per application: policy / wallet UX & cross-application RU; \num{15} owners hold \num{223} applications \\
Q16 & A, C & which wallet implements $\zaddr\to\taddr$ with \texttt{OP\_RETURN} & determines who can actually use Construction~\ref{con:pp-a} \\
Q17 & A & proof-count cap / \texttt{librustzcash} upgrade / none & Property~\ref{prop:pp-verify} \\
\bottomrule
\end{tabular}
\end{table}

\begin{remark}[Every question in this section is conditional on the pool surviving]
\label{rem:pp-conditional}
Of Table~\ref{tab:pp-open}'s seventeen rows, Q1 is withdrawn and Q2 answered
(\Cref{rem:q1q2-resolved}), Q3 is settled (\S7, O3) and Q4 is mooted by retirement; the rest remain
open decisions in form, and Q6, Q8, Q12 and Q17 are called out separately above because they gate
constructibility, cost or a security property. No
value in this section was chosen by this drafter, and that remains deliberate --- but the reason has
changed, and the change is worth stating rather than leaving implicit.

\Cref{sec:shielded-retirement} records the option the team has since taken: retiring both shielded
pools outright, which the measurements there favour \srcintent{08-answers-round-2.md, round 10}.
\textbf{With that option taken, this entire section is moot} --- not deferred, but inapplicable,
because there is no pool to pay into, no proving system to choose, no denomination grid to design
and no turnstile question beyond the drain window. Q1's ``reopening is decided, only the height is
open'' framing was recorded before that option was on the table and is now withdrawn.

The paper therefore does not fill these parameters, and would not be improved by filling them: every
value would be conditional on a pool that will not exist. What this section provides regardless is
the requirement set --- what a
private payment for a public obligation must satisfy --- and that is independent of whether Flux
carries a shielded pool. The construction is what depends on the answer.
\end{remark}

Four measurements are also outstanding and are not decisions: \textbf{M1}, the Sapling pool in notes
and holders (Open Problem~\ref{op:pp-m1}); \textbf{M2}, renewal churn, which converts
\S\ref{sec:pp-anonymity}'s upper-bound rate into a rate; \textbf{M3}, CumulusVPN's payment rate, which
is the anonymity set for Construction~\ref{con:pp-b}; and \textbf{M4}, $t_{G16}$ on tier hardware,
which converts Property~\ref{prop:pp-verify} from a parameter to a number. \textbf{M5}, whether the
daemon-proxy address index returns empty-\texttt{vin} Sapling transactions paying $A_{\mathrm{sink}}$,
is resolved, favourably: it does (\S\ref{sec:pp-deployed}, interface I1), so it gates
Construction~\ref{con:pp-a} only on reopening (P1) and no longer stands as an open test.

\subsection{Comparison with prior art}
\label{sec:pp-related}

\begin{table}[H]
\centering
\scriptsize
\caption{Prior art against the requirement pair of \S\ref{sec:pp-requirements}. The column that
separates the field is \emph{where soundness is anchored}: on a chain, or at an issuer.}
\label{tab:pp-related}
\begin{tabular}{@{}lp{0.20\linewidth}p{0.22\linewidth}lp{0.22\linewidth}@{}}
\toprule
\textbf{System} & \textbf{What it hides} & \textbf{How the obligation is verified} & \textbf{Soundness} & \textbf{What Flux does differently} \\
\midrule
Zerocash / Zcash Sapling \citep{bensasson2014zerocash,zcashprotocol} &
sender, receiver and amount for $\zaddr\to\zaddr$; sender for $\zaddr\to\taddr$ &
no notion of an obligation beyond the transfer; no recurring-payment primitive &
chain &
attaches a \emph{public} obligation $H_\app$ to a $\zaddr\to\taddr$ and lets \num{6489} independent
verifiers act on it; the shielded-subscription problem Zcash never solved is shape C, and the answer
here is child-key delegation, not a protocol primitive \\
\addlinespace
Groth16 and its ceremony \citep{groth2016size,bowe2017mpc} &
--- &
--- &
--- &
Flux inherits the 2018 Zcash Sapling parameters verbatim and has run no ceremony of its own; every
\emph{new} statement would need one (Q8) \\
\addlinespace
Halo, no trusted setup \citep{bowe2019halo} &
--- &
--- &
--- &
removes the inherited ceremony, which directly serves \S20's decentralization argument; it is
\textbf{not} post-quantum and the two claims are kept separate \\
\addlinespace
Fixed-denomination mixers &
the link deposit\,$\leftrightarrow$\,withdrawal within a denomination &
on-chain nullifier set plus Merkle membership &
chain &
the identical trade-off, measured here on Flux's own demand (Table~\ref{tab:pp-denom}); shape A needs
no separate pool because Sapling \emph{is} one \\
\addlinespace
Semaphore-style membership signalling &
which member signalled &
Merkle membership plus a per-external-nullifier nullifier &
chain or verifier-side &
Construction~\ref{con:pp-b} is this shape with the membership set derived from \emph{payments on
chain}, so the set needs no admission authority \\
\addlinespace
Lightning \citep{poon2016lightning} &
the payer from third parties (onion routing); the payment from non-participants &
the recipient alone; the preimage is a receipt &
channel counterparties &
Lightning has no \emph{public} obligation: a receipt convinces one party. Flux needs \num{6489}
independent verifiers holding no secret, which is precisely what point-to-point payment privacy
cannot provide \\
\addlinespace
Chaumian blind signatures; Privacy-Pass-style tokens &
the link issuance\,$\leftrightarrow$\,redemption &
the issuer's public key, with a per-verifier or shared spent set &
\textbf{issuer} &
admissible under V0, which constrains \emph{verification} only --- but soundness leaves the chain
(Remark~\ref{rem:pp-voucher}); Construction~\ref{con:pp-b} replaces the issuer with the chain and pays
a circuit for it \\
\addlinespace
Anonymous credentials (CL-style; algebraic-MAC keyed-verification) &
multi-show unlinkability with provable attributes such as \texttt{paid\_until} &
the issuer's key; keyed-verification requires the \emph{verifier} to hold the issuer secret &
issuer &
keyed-verification violates V0 outright for \num{6489} independent verifiers --- fine for a
single-gateway model, wrong for this node model \\
\addlinespace
Compact e-cash / rechargeable anonymous wallets &
which coin, which recharge &
bank-issued coins with double-spend detection &
the bank &
the closest prior art for shape C: the shielded notes are the coins and the chain is the bank, but
delegation to a renewer (R2) is where the analogy stops \\
\bottomrule
\end{tabular}
\end{table}

The prior art named in the last four rows of Table~\ref{tab:pp-related} is
\citet{chaum1982blind} for blind signatures, the Privacy Pass line of anonymous authorization tokens
\citep{privacypass}, \citet{camenisch2001credentials} for anonymous credentials and the
keyed-verification line that follows it, and \citet{camenisch2005ecash} for compact e-cash.

The combination that is specific to Flux is worth stating in one sentence, because it is what makes
this section a contribution rather than a survey: \textbf{a public obligation that must be checkable
by thousands of independent verifiers holding no shared secret, paid out of a shielded pool that
already exists but cannot be entered, at a demand level where the anonymity set rather than the
cryptography is the binding constraint.} Two of the three difficulties are not cryptographic. The
construction for the common case is short and uses circuits that have been running on this network
since 2019; the honesty about the set, and about the consensus rule that closed the pool in 2021, is
the part that took work.

%% Drain this section's float queue before the next begins.
\FloatBarrier
